\documentclass[../main.tex]{subfiles}
\begin{document}

\chapter{VLIW}
\lessoninfo{
  video={Lesson 11，動画 1--9},
  textbook={H\&P \cite{hennessy2017} §3.7，付録 H},
  keywords={VLIW，アウトオブオーダ・スーパースカラ，インオーダ・スーパースカラ，NOP，後方互換性，不規則なコード，コードの肥大化，完全プレディケーション，分岐ヒント，ストップビット，Itanium，EPIC，DSP},
}

第10章では，コンパイラがプログラムを書き換えることで，プロセッサが
1 サイクルに完了する命令の数を増やせることを示した．コンパイラは依存の
連鎖を短くし，独立な命令を隣り合わせに配置し，分岐と呼び出しを除去する．
VLIW（very long instruction word）プロセッサは，この考え方を突き詰めた
ものである．どの操作を並列に実行するかをコンパイラだけが決め，その決定を
命令に符号化するので，ハードウェアには第7〜9章のアウトオブオーダ
プロセッサが持つ依存追跡の機構が一切必要ない．この章では，VLIW を
スーパースカラの設計と比較し，その利点と欠点を検討し，この手法が失敗した
分野と成功した分野を示す．

\section{スーパースカラプロセッサと VLIW プロセッサ}
\label{sec:l11-superscalar-vs-vliw}

スーパースカラ\index{すーぱーすから@スーパースカラ}プロセッサ（第8章）は，
1 サイクルに最大 $N$ 個の命令を実行する．\source{Lesson 11，動画 1--3；
H\&P §3.7；\cite{fisher1983}．}設計による違いは，一緒に実行できる $N$ 個の
命令を見つけるためにいくつの命令を調べるか，そしてその探索のどれだけを
コンパイラに任せるかにある（\cref{fig:l11-windows,tab:l11-compare}）．
アウトオブオーダの設計は，リネーミング，ウィンドウ内のスケジューリング，
リオーダバッファのための複雑なハードウェアと引き換えに，どのような
プログラムの並列性も自力で見つける（第6〜8章）．インオーダの設計は，他の
命令を待たなければならない最初の命令で止まるので，独立な命令を隣り合わせに
配置するコンパイラの助けを必要とする（第10章）．VLIW プロセッサは，
コンパイラが互いに独立にした最大 $N$ 個の操作からなる 1 つの命令を，
何の検査もせずに実行する．

\begin{figure}[htbp]
  \centering
  % What each design examines to find N = 4 operations for one cycle:
% out-of-order superscalar (window >> N), in-order superscalar (next N),
% VLIW (one long instruction).  Shaded = executed in this cycle.
\begin{tikzpicture}[x=1cm,y=1cm, font=\footnotesize\sffamily, >={Stealth[length=1.6mm]},
  ins/.style={draw, minimum width=0.55cm, minimum height=0.45cm, inner sep=0pt, fill=white},
  run/.style={ins, fill=ln-accent!20},
  lab/.style={anchor=east, align=right, text width=2.6cm},
  win/.style={ln-accent, line width=0.6pt}]
  % --- row labels
  \node[lab] at (2.9,0)    {\iflangja{アウトオブオーダ\\スーパースカラ}{out-of-order\\superscalar}};
  \node[lab] at (2.9,-1.4) {\iflangja{インオーダ\\スーパースカラ}{in-order\\superscalar}};
  \node[lab] at (2.9,-3.4) {VLIW};
  % --- program order arrow
  \draw[->, ln-muted] (3.2,-4.15) -- (6.2,-4.15)
    node[right, font=\scriptsize\sffamily] {\iflangja{プログラム順}{program order}};
  % --- row 1: out-of-order, window of 12, four scattered picks
  \foreach \i in {0,...,11} {
    \node[ins] (o\i) at (3.475+0.55*\i,0) {};
  }
  \foreach \i in {0,1,4,8} { \node[run] at (o\i) {}; }
  \node[ln-muted] at (10.25,0) {\dots};
  \draw[win] (3.2,0.35) -- (3.2,0.45) -- (9.8,0.45) -- (9.8,0.35);
  % --- row 2: in-order, window of 4, fourth depends on second
  \foreach \i in {0,...,11} {
    \node[ins] (n\i) at (3.475+0.55*\i,-1.4) {};
  }
  \foreach \i in {0,1,2} { \node[run] at (n\i) {}; }
  \node[ln-muted] at (10.25,-1.4) {\dots};
  \draw[win] (3.2,-1.05) -- (3.2,-0.95) -- (5.4,-0.95) -- (5.4,-1.05);
  \draw[->, ln-caution] (n1.south) to[bend right=50]
    node[below, font=\scriptsize\sffamily] {\iflangja{依存}{dependence}} (n3.south);
  % --- row 3: VLIW, long instructions of four slots
  \foreach \j in {0,1,2} {
    \draw[fill=white] (3.2+2.2*\j,-3.625) rectangle (5.4+2.2*\j,-3.175);
    \foreach \k in {1,2,3} {
      \draw[ln-rule] (3.2+2.2*\j+0.55*\k,-3.625) -- (3.2+2.2*\j+0.55*\k,-3.175);
    }
  }
  \fill[ln-accent!20] (3.2,-3.625) rectangle (5.4,-3.175);
  \draw (3.2,-3.625) rectangle (5.4,-3.175);
  \foreach \k in {1,2,3} { \draw (3.2+0.55*\k,-3.625) -- (3.2+0.55*\k,-3.175); }
  \draw[line width=0.8pt] (5.4,-3.625) -- (5.4,-3.175) (7.6,-3.625) -- (7.6,-3.175);
  \node[ln-muted] at (10.25,-3.4) {\dots};
  \draw[win] (3.2,-3.05) -- (3.2,-2.95) -- (5.4,-2.95) -- (5.4,-3.05);
  % --- window labels
  \node[font=\scriptsize\sffamily, text=ln-accent, anchor=south] at (6.5,0.42)
    {\iflangja{調べる範囲：$N$ よりはるかに多い命令}{examined: $\gg N$ instructions}};
  \node[font=\scriptsize\sffamily, text=ln-accent, anchor=south west, inner xsep=0pt] at (3.2,-0.95)
    {\iflangja{調べる範囲：次の $N$ 命令}{examined: the next $N$ instructions}};
  \node[font=\scriptsize\sffamily, text=ln-accent, anchor=south west, inner xsep=0pt] at (3.2,-2.95)
    {\iflangja{調べる範囲：1 命令（$N$ 個の操作）}{examined: 1 instruction ($N$ operations)}};
\end{tikzpicture}

  \caption{1 サイクルを最大 $N=4$ 個の操作で埋めるために，各設計が調べる
    範囲．網掛けの箱がこのサイクルで実行される．インオーダの設計は，
    グループ内の先行する命令に依存する最初の命令で止まる．}
  \label{fig:l11-windows}
\end{figure}

\begin{definition}[VLIW]
  \label{def:l11-vliw}
  \term{VLIW}[very long instruction word]（very long instruction word）
  アーキテクチャとは，各命令が最大 $N$ 個の操作からなり，それらの操作が
  同じサイクルで並列に実行される命令セットアーキテクチャである．各操作は
  通常の ISA の 1 命令に相当する．1 つの命令を構成する操作を選び，それらが
  独立であることを保証するのはコンパイラの役割であり，プロセッサは依存の
  検査を一切行わない．
\end{definition}

$N$ 個の操作を収める命令は通常の命令のおよそ $N$ 倍の長さになり，これが
このアーキテクチャの名前の由来である．1 命令が最大 $N$ 個の通常の命令の
仕事をするので，1 サイクルに 1 命令を完了する VLIW プロセッサは，IPC が
最大~$N$ のスーパースカラプロセッサに相当する（第6章）．\caution{VLIW
プログラムの命令数と CPI は，通常の ISA のものとは比較できない．サイクル
当たりの操作数で比較すること．}

\begin{table}[htbp]
  \centering\small
  \caption{1 サイクルに最大 $N$ 個の操作を実行するスーパースカラプロセッサ
    と VLIW プロセッサ．}
  \label{tab:l11-compare}
  \begin{tabular}{@{}lccc@{}}
    \toprule
                         & アウトオブオーダ  & インオーダ      & VLIW \\
    \midrule
    サイクル当たりの実行 & $\le N$ 命令      & $\le N$ 命令    & 1 命令（$\le N$ 操作） \\
    調べる範囲           & $\gg N$ 命令      & 次の $N$ 命令   & 1 命令 \\
    ハードウェア         & 複雑              & より単純        & 最も単純 \\
    コンパイラ           & 助けになりうる    & 助けが必須      & すべてを担う \\
    \bottomrule
  \end{tabular}
\end{table}

\begin{example}[命令への操作の詰め込み]
  \label{ex:l11-packing}
  次の断片は配列の 4 要素を加算し，その和をストアする．コメントは RAW 依存
  （第3章）を示す．4 つのロードは R8 を読むだけであり，互いに独立である．
  \begin{lstlisting}[language={[x86masm]Assembler}, numbers=none]
I1: LW  R1, 0(R8)
I2: LW  R2, 4(R8)
I3: LW  R3, 8(R8)
I4: LW  R4, 12(R8)
I5: ADD R5, R1, R2    ; I1，I2 と RAW
I6: ADD R6, R3, R4    ; I3，I4 と RAW
I7: ADD R7, R5, R6    ; I5，I6 と RAW
I8: SW  R7, 16(R8)    ; I7 と RAW
\end{lstlisting}
  VLIW コンパイラは操作をプログラム順に取り出し，それぞれを，依存する
  すべての操作より後にあり，かつ空きスロットが残っている最も早い命令に
  配置する．余ったスロットは NOP で埋める（\cref{tab:l11-packing}）．
  通常の ISA なら 8 命令が必要なところ，スロットが 4 つなら 8 個の操作は
  4 命令，すなわち 4 サイクルに収まり，スロットが 2 つなら 5 命令を占める．
\end{example}

\begin{table}[htbp]
  \centering\small
  \caption{\cref{ex:l11-packing} の断片を，操作スロットが 4 つの VLIW ISA
    と 2 つの VLIW ISA に向けてコンパイルした結果．各行が 1 命令である．
    サイズは 32 ビットのスロットを仮定している．}
  \label{tab:l11-packing}
  \begin{tabular}{@{}cll@{}}
    \toprule
    命令         & 4 スロット         & 2 スロット \\
    \midrule
    1            & I1, I2, I3, I4     & I1, I2 \\
    2            & I5, I6, NOP, NOP   & I3, I4 \\
    3            & I7, NOP, NOP, NOP  & I5, I6 \\
    4            & I8, NOP, NOP, NOP  & I7, NOP \\
    5            & ---                & I8, NOP \\
    \midrule
    NOP スロット & 16 個中 8 個       & 10 個中 2 個 \\
    コードサイズ & 512~ビット         & 320~ビット \\
    \bottomrule
  \end{tabular}
\end{table}

\section{利点と欠点}
\label{sec:l11-good-bad}

並列性の探索をコンパイラに任せることには 3 つの利点がある．
\source{Lesson 11，動画 4--5；H\&P §3.7，付録 H．}
\begin{itemize}
  \item \textbf{コンパイラには時間が十分にある．}アウトオブオーダ
    プロセッサは，毎クロックサイクルのごく一部の時間で，しかも自身の
    ウィンドウの中だけで独立な命令を見つけなければならない．コンパイラは
    一度だけ実行され，プログラム全体を解析できる．
  \item \textbf{ハードウェアが単純になる．}リネーミング，ウィンドウ内の
    スケジューリング，リオーダバッファがないので，VLIW プロセッサは同じ
    サイクル当たり操作数を実現するのに必要な面積と電力が少なくて済み，
    エネルギー効率を高くできる．
  \item \textbf{ループと規則的なコードに適する．}配列をループで処理する
    場合，依存は少なく，コンパイル時に把握できる．ループ展開
    \index{るーぷてんかい@ループ展開}（第10章）が反復間の分岐を除去する
    ので，独立な操作がスロットを埋める．
\end{itemize}
一方で，欠点は 4 つある．
\begin{itemize}
  \item \textbf{レイテンシがコンパイル時に分からない．}コンパイラは，各操作
    がそのサイクル内に完了するものとして操作をグループ化する．ある操作が
    1 サイクルで終わらないと---例えばキャッシュミス（第12章）を起こした
    ロード---後続のすべての操作が，独立なものまで含めてそれを待つ．一方，
    アウトオブオーダプロセッサはそれに依存しない命令の実行を続ける．
  \item \textbf{不規則なコードは難しい．}コンパイラが見るのはプログラム
    であって，その実行ではない．2 つの操作が独立であることを証明できない
    とき---同じ位置を指しうるポインタを介したストアとロードや，データに
    依存する分岐の両側にある操作---コンパイラはそれらを別々の命令に置か
    なければならない．一方，アウトオブオーダプロセッサは実行時に実際の
    アドレスを知り，分岐を予測する（第4章と第9章）．
  \item \textbf{コードの肥大化．}独立な操作が $N$ 個未満しかないたびに
    NOP がスロットを埋め，プログラムが大きくなる．\cref{tab:l11-packing}
    の 4 スロット版は半分が NOP で 512~ビットを占めるが，通常の 32 ビット
    命令 8 個なら 256~ビットで済む．命令の圧縮（\cref{sec:l11-instructions}）
    はこのオーバーヘッドを減らす．
  \item \textbf{新しいプロセッサには再コンパイルしたプログラムが必要である．}%
    操作のグループ化とスロット数はコンパイル時に固定されるので，後方互換性
    （後述）が損なわれる．
\end{itemize}

\begin{definition}[後方互換性]
  プロセッサが\term[こうほうごかんせい]{後方互換性}[backward compatibility]%
  を持つとは，同じファミリの以前のプロセッサ向けにコンパイルされた
  プログラムを，変更せずに実行できることである．
\end{definition}

スーパースカラプロセッサでは，後方互換性に性能向上が伴う．幅の広い
アウトオブオーダプロセッサは，\emph{古い}バイナリからも再コンパイルなしで
より多くの並列性を見つける．スロット数を増やした新しい VLIW プロセッサは，
命令形式が変わったために古い VLIW バイナリをまったく実行できないか，
実行できても，そのバイナリのコンパイル対象だったプロセッサよりサイクル
当たりの操作数が増えない．

\begin{example}[幅の広い VLIW プロセッサ上の古いバイナリ]
  \label{ex:l11-compat}
  2 スロット形式も受け付ける 4 スロットのプロセッサは，
  \cref{ex:l11-packing} の 2 スロット版を，古いプロセッサと同じく
  5 サイクルで実行する．4 サイクルで済むのは再コンパイルした 4 スロット版
  だけである．8 個の通常の命令は，同じ数え方をすると，再コンパイルなしで，
  幅 2 のアウトオブオーダプロセッサでは 5 サイクル，幅 4 のものでは
  4 サイクルかかる．
\end{example}

\begin{keypoint}
  VLIW は，ハードウェアの複雑さをコンパイラへの依存と引き換えにする．
  コードが規則的で，まさにそのプロセッサ向けにコンパイルされている場合には
  見合うが，並列性や操作のレイテンシが実行時にしか分からない場合や，古い
  バイナリを新しいプロセッサでより速く実行しなければならない場合には不利に
  なる．
\end{keypoint}

\section{VLIW 命令}
\label{sec:l11-instructions}

VLIW 命令は $N$ 個の操作スロットの並びである（\cref{fig:l11-format}）．
\source{Lesson 11，動画 6；H\&P 付録 H．}各スロットは通常の ISA の命令と
同じように符号化され，通常の命令コード全体から選ばれた固有の命令コードと，
固有のオペランドを持つ．

\begin{figure}[htbp]
  \centering
  % A VLIW instruction with N = 4 operation slots, each encoded like an
% ordinary instruction, decoded and executed in parallel.
\begin{tikzpicture}[x=1cm,y=1cm, font=\footnotesize\sffamily, >={Stealth[length=1.6mm]},
  unit/.style={draw, fill=ln-box, minimum width=1.9cm, minimum height=0.5cm, inner sep=1pt}]
  % --- the instruction word: four contiguous slots of 2.45 cm
  \foreach \s in {0,...,3} {
    \pgfmathsetmacro{\x}{2.45*\s}
    \fill[ln-accent!15] (\x,0) rectangle (\x+0.95,0.55);
    \foreach \d in {0.95,1.45,1.95} { \draw[ln-rule] (\x+\d,0) -- (\x+\d,0.55); }
    \node[font=\scriptsize\sffamily] at (\x+0.475,0.275) {\iflangja{{\tiny 命令コード}}{opcode}};
    \node[font=\scriptsize\sffamily] at (\x+1.2,0.275) {rd};
    \node[font=\scriptsize\sffamily] at (\x+1.7,0.275) {rs};
    \node[font=\scriptsize\sffamily] at (\x+2.2,0.275) {rt};
    \pgfmathtruncatemacro{\n}{\s+1}
    \node[font=\scriptsize\sffamily, text=ln-muted, anchor=south] at (\x+1.225,0.55)
      {\iflangja{スロット \n}{slot \n}};
    % decode and execute
    \node[unit] (d\s) at (\x+1.225,-0.9) {\iflangja{デコード}{decode}};
    \node[unit] (e\s) at (\x+1.225,-1.8) {\iflangja{実行}{execute}};
    \draw[->] (\x+1.225,0) -- (d\s);
    \draw[->] (d\s) -- (e\s);
    \draw[->] (e\s) -- (\x+1.225,-2.6);
  }
  \foreach \s in {1,2,3} { \draw[line width=0.9pt] (2.45*\s,0) -- (2.45*\s,0.55); }
  \draw[line width=0.9pt] (0,0) rectangle (9.8,0.55);
  % --- architectural state
  \draw[fill=white] (0,-2.6) rectangle (9.8,-3.1);
  \node at (4.9,-2.85) {\iflangja{レジスタとメモリ}{registers and memory}};
  % --- marker for the whole word
  \draw[ln-accent, line width=0.6pt] (0,1.0) -- (0,1.1) -- (9.8,1.1) -- (9.8,1.0);
  \node[font=\scriptsize\sffamily, text=ln-accent, anchor=south] at (4.9,1.1)
    {\iflangja{1 命令 $=$ $N$ 個の独立な操作}{one instruction $=$ $N$ independent operations}};
\end{tikzpicture}

  \caption{4 つの操作スロットを持つ VLIW 命令．各スロットは他のスロットと
    独立にデコードされ，実行される．スロット間の依存を検査するハードウェア
    はない．}
  \label{fig:l11-format}
\end{figure}

1 つの命令に含まれる操作は独立でなければならない．どの操作も他の操作の
結果を使ってはならず，2 つの操作が同じ位置に書き込んではならない．
ハードウェアはこれを検査しないので，正しくグループ化する責任はコンパイラ
にある．\emph{連続する}命令の操作どうしは依存してよい．
\cref{tab:l11-packing} で I5 が I1 に依存しているのがその例である．
\caution{独立性が要求されるのは 1 つの命令の中だけであり，命令と命令の
間では要求されない．}

コンパイラは，\cref{ex:l11-packing} のように命令スケジューリングによって
並列性を引き出す．第10章の他の技法もこれを助ける．木の高さ削減
\index{きのたかささくげん@木の高さ削減}は依存の連鎖を短くして，より多くの
操作が 1 つの命令を共有できるようにし，ループ展開とインライン展開
\index{いんらいんてんかい@インライン展開}は，操作を隔ててしまう分岐，
呼び出し，リターンを除去する．VLIW プロセッサにはリネーミングの
ハードウェアがないので，偽の依存（第3章）はコンパイラ自身が，他で使われて
いないレジスタに結果を書き込むことで除去する．そのためのレジスタを，命令
セットは数多く用意している．

分岐の後の操作は分岐の一方の経路でしか実行されないので，分岐はグループ化
を制限する．そのため VLIW の命令セットはプレディケーション（第5章），
理想的には完全プレディケーション
\index{かんぜんぷれでぃけーしょん@完全プレディケーション}を十分に
サポートし，コンパイラが分岐を if 変換\index{if へんかん@if 変換}して
両方の経路の操作をグループ化できるようにする．残った分岐には
\term[ぶんきひんと]{分岐ヒント}[branch hint]を付けることができ，
コンパイラはこれによって，分岐がどちらに進むと予想しているかをプロセッサ
に伝える．

固定スロットは，独立な操作が $N$ 個未満しかないたびに領域を無駄にする．
圧縮された命令は操作を NOP なしで順に並べて格納し，一緒に実行される各
グループの最後の操作に\term[すとっぷびっと]{ストップビット}[stop bit]で
印を付ける．\cref{tab:l11-packing} の 4 スロット版のプログラムでは I4，
I6，I7，I8 がストップビットを持ち，操作当たり 33~ビットとすると，コードは
512~ビットではなく 264~ビットを占める．

\section{実際の VLIW プロセッサ}
\label{sec:l11-examples}

VLIW の手法がどこで失敗し，どこで成功するかを，2 種類のプロセッサが
示している．\source{Lesson 11，動画 7--9；H\&P §3.7，付録 H．}

\subsection{Itanium}
\label{sec:l11-itanium}

VLIW 風の命令を汎用コンピューティングに持ち込もうとした最も有名な試みは，
Intel が IA-64 命令セットのために開発したプロセッサファミリ
\term{Itanium}[Itanium]である．設計者たちはこの手法を EPIC（explicitly
parallel instruction computing）と呼んだ \cite{schlansker2000}．
\margin{IA-64 も同じようにストップを用いて，独立な命令のグループの終わり
を示す（\cref{sec:l11-instructions}）．}しかし，この命令セットには多くの
追加機能が積み重なり，そのハードウェアが非常に複雑になったため，VLIW の
主な利点が失われた．不規則なコードからはやはりコンパイラがほとんど並列性を
見つけられず，Itanium がアウトオブオーダプロセッサに取って代わることは
なかった．

\subsection{DSP}
\label{sec:l11-dsp}

\term{DSP}[digital signal processor]（digital signal processor）とは，
音声，映像，無線などのデジタル信号の処理に特化したプロセッサである．
この市場は VLIW の長所（\cref{sec:l11-good-bad}）を生かせる．
\margin{H\&P は VLIW の手法の例として TI の C6x 系 DSP を挙げている．}
\begin{itemize}
  \item 同じ計算を多数のサンプルに適用するループがコードの大半を占める．
    コンパイラが並列性を十分に見つけられる，規則的なコードである．
  \item プログラムは通常，それが動作する機器の特定の DSP 向けに
    コンパイルされるので，バイナリの後方互換性はほとんど問題にならない．
  \item DSP は電力の限られた機器に搭載されるので，エネルギー効率の高い
    ハードウェアが効果を発揮する．
\end{itemize}

\needspace{6\baselineskip}
\begin{keypoint}[章のまとめ]
  アウトオブオーダのスーパースカラプロセッサは $N$ よりはるかに多い命令を
  調べ，コンパイラの助けを必要としない．インオーダのものは $N$ 個の命令を
  調べ，その助けを必要とする．VLIW プロセッサは最大 $N$ 個の独立な操作から
  なる 1 命令を調べ，全面的にコンパイラに頼る．VLIW のハードウェアは単純で
  エネルギー効率が高く，規則的なループのコードで良い性能を示す．しかし
  新しいプロセッサごとにバイナリを再コンパイルしなければならず，不規則な
  コードからはほとんど並列性が得られず，1 つのロードのキャッシュミスが後続の
  すべての操作を遅らせ，ストップビットで圧縮しない限り NOP がコードを
  肥大化させる．Itanium は汎用コンピューティングで失敗し，VLIW が適するのは
  DSP である．第12章では，サイクル当たり命令数からメモリへと話題を移す．
\end{keypoint}

\end{document}
